% ====================================================================
% PREÂMBULO UNIVERSAL LATEX (TEMPLATE DE INCLUDES SEGURO E COMPLETO)
% ====================================================================
% Como usar: Salve este arquivo como 'includes.tex' no mesmo diretório
% do seu documento principal e adicione % ====================================================================
% PREÂMBULO UNIVERSAL LATEX (TEMPLATE DE INCLUDES SEGURO E COMPLETO)
% ====================================================================
% Como usar: Salve este arquivo como 'includes.tex' no mesmo diretório
% do seu documento principal e adicione % ====================================================================
% PREÂMBULO UNIVERSAL LATEX (TEMPLATE DE INCLUDES SEGURO E COMPLETO)
% ====================================================================
% Como usar: Salve este arquivo como 'includes.tex' no mesmo diretório
% do seu documento principal e adicione % ====================================================================
% PREÂMBULO UNIVERSAL LATEX (TEMPLATE DE INCLUDES SEGURO E COMPLETO)
% ====================================================================
% Como usar: Salve este arquivo como 'includes.tex' no mesmo diretório
% do seu documento principal e adicione \input{includes} antes de
% \begin{document}.
% ====================================================================

% --- 0. COMPILAÇÃO E FONTE ---
\usepackage[T1]{fontenc}                    % Codificação de fonte necessária para o pdfLaTeX
\usepackage[utf8]{inputenc}                 % Garante o suporte a acentuação nativa UTF-8
\usepackage{helvet}                         % Carrega a família Helvetica (equivalente visual ao Arial)
\renewcommand{\familydefault}{\sfdefault}   % Define a fonte sem serifa (Helvetica) como padrão do documento

% --- 1. IDIOMA E AJUSTES TIPOGRÁFICOS ---
\usepackage[brazil]{babel}          % Tradução de termos padrão (Figuras, Tabelas, Sumário)
\usepackage{microtype}              % Melhorias tipográficas estruturais (ajuste de espaçamento)

% --- 2. MATEMÁTICA AVANÇADA ---
\usepackage{amsmath}                    % Comandos matemáticos básicos da AMS
\usepackage{amsfonts}                   % Fontes matemáticas adicionais
\usepackage{amssymb}                    % Símbolos matemáticos diversos
\usepackage{mathtools}                  % Extensões e correções para o amsmath
\usepackage{bm}                         % Negrito para símbolos matemáticos (vetores, matrizes)
\usepackage{siunitx}                    % Formatação padronizada de unidades do SI (crucial para Engenharia)
\sisetup{output-decimal-marker = {,}}   % Configura a vírgula como separador decimal para PT-BR/PT-PT

% --- 3. LAYOUT, PÁGINA E ESPAÇAMENTO ---
\usepackage{geometry}       % Ajuste flexível das margens da página
\geometry{
    a4paper,
    top=30mm,
    bottom=20mm,
    left=30mm,
    right=20mm,
    % includehead,
    % includefoot,
    headheight=1.4cm,
    headsep=0.4cm,
    % headsep=0.6cm,
    footskip=1cm
    % footskip=1.4cm
}

\usepackage{fancyhdr}               % Customização avançada de cabeçalhos e rodapés
\usepackage{lastpage}               % Permite referenciar o número da última página (ex: "2 de 5")
\usepackage{setspace}               % Controle do espaçamento entre linhas (\onehalfspacing, etc.)
\usepackage{multicol}               % Suporte para layouts de múltiplas colunas na mesma página
\usepackage{indentfirst}            % Garante que o primeiro parágrafo de cada seção seja indentado

% --- 4. ELEMENTOS GRÁFICOS E CORES ---
\usepackage{graphicx}               % Inclusão de imagens e gráficos externos
\usepackage{transparent}            % Permite definir opacidade de imagens (ex.: logos com 70% de opacidade)
\usepackage{xcolor}                 % Suporte avançado e paletes de cores customizadas
\usepackage{tikz}                   % Ferramenta poderosa para criação de diagramas nativos
\usetikzlibrary{calc}               % Permite aritmética de coordenadas (usada na faixa do cabeçalho)
\usepackage{pgfplots}               % Geração de gráficos matemáticos de alta precisão
\pgfplotsset{compat=1.18}           % Garante compatibilidade do pgfplots

% --- 5. TABELAS PROFISSIONAIS ---
\usepackage{booktabs}               % Linhas de tabela padrão publicação (\toprule, \midrule, \bottomrule)
\usepackage{array}                  % Extensões para propriedades das colunas em tabelas
\usepackage{colortbl}               % Permite colorir linhas/colunas/células de tabelas (\rowcolor, \cellcolor)
\usepackage{multirow}               % Permite mesclar linhas em tabelas
\usepackage{tabularx}               % Tabelas com largura total fixa e colunas auto-ajustáveis (X)
\usepackage{longtable}              % Suporte para tabelas que quebram ao longo de múltiplas páginas

% --- 6. FIGURAS, LEGENDAS E FLUTUANTES ---
\usepackage{float}                  % Controle estrito de posicionamento (modificador [H])
\usepackage{caption}                % Customização estética das legendas de figuras e tabelas
\usepackage{subcaption}             % Criação de subfiguras e sub-legendas organizadas

% Legendas de figuras/tabelas usam o tamanho "menor" do padrão (Arial 10)
\DeclareCaptionFont{arialdez}{\fontsize{10}{12}\selectfont}
\captionsetup{font=arialdez, labelfont={arialdez,bf}}

% --- 7. LISTAS E ENUMERAÇÕES ---
\usepackage{enumitem}               % Customização total de listas (itemize, enumerate, description)

% --- 8. CÓDIGO-FONTE E CAIXAS DE DESTAQUE ---
\usepackage{listings}               % Inclusão e estilização de blocos de código-fonte
\usepackage{tcolorbox}              % Criação de caixas de texto coloridas e estilizadas de alta qualidade
\tcbuselibrary{skins,breakable}     % Permite caixas que quebram páginas e estilos complexos

% --- 9. TÍTULOS E SÍMBOLOS DIVERSOS ---
\usepackage{titlesec}               % Customização visual de \section, \subsection etc.
\usepackage{pifont}                 % Símbolos diversos (ex: \ding{51} para marcações de check)

% --- 10. REFERÊNCIAS E HIPERLINKS (Devem ser carregados por último) ---
\usepackage{url}                    % Formatação correta de URLs
\usepackage{hyperref}               % Transforma referências e sumários em links clicáveis interativos
\hypersetup{
    colorlinks=true,                % True para links coloridos, False para caixas coloridas
    linkcolor=black,                % Cor dos links internos (equações, seções)
    citecolor=black,                % Cor das citações bibliográficas
    urlcolor=black                  % Cor de links externos da internet
}
\usepackage{cleveref}               % Referências cruzadas inteligentes (escreve "Figura 1" automaticamente)

% Configuração de tradução do cleveref para português (garantia adicional)
\crefname{figure}{figura}{figuras}
\Crefname{figure}{Figura}{Figuras}
\crefname{table}{tabela}{tabelas}
\Crefname{table}{Tabela}{Tabelas}
\crefname{equation}{equação}{equações}
\Crefname{equation}{Equação}{Equações}

% --------------------------------------------------------------
% Aviso "Token not allowed in a PDF string (Unicode): removing
% `subscript'" do hyperref: aparece sempre que um campo de formulário
% (\TextField/\CheckBox, usados em campos-formulario.tex) é criado
% dentro de uma macro — é um comportamento interno conhecido do
% hyperref com AcroForm, inofensivo (não afeta o PDF gerado). A linha
% abaixo apenas filtra esse aviso específico da saída do compilador.
\usepackage{silence}
\WarningFilter{hyperref}{Token not allowed in a PDF string} antes de
% \begin{document}.
% ====================================================================

% --- 0. COMPILAÇÃO E FONTE ---
\usepackage[T1]{fontenc}                    % Codificação de fonte necessária para o pdfLaTeX
\usepackage[utf8]{inputenc}                 % Garante o suporte a acentuação nativa UTF-8
\usepackage{helvet}                         % Carrega a família Helvetica (equivalente visual ao Arial)
\renewcommand{\familydefault}{\sfdefault}   % Define a fonte sem serifa (Helvetica) como padrão do documento

% --- 1. IDIOMA E AJUSTES TIPOGRÁFICOS ---
\usepackage[brazil]{babel}          % Tradução de termos padrão (Figuras, Tabelas, Sumário)
\usepackage{microtype}              % Melhorias tipográficas estruturais (ajuste de espaçamento)

% --- 2. MATEMÁTICA AVANÇADA ---
\usepackage{amsmath}                    % Comandos matemáticos básicos da AMS
\usepackage{amsfonts}                   % Fontes matemáticas adicionais
\usepackage{amssymb}                    % Símbolos matemáticos diversos
\usepackage{mathtools}                  % Extensões e correções para o amsmath
\usepackage{bm}                         % Negrito para símbolos matemáticos (vetores, matrizes)
\usepackage{siunitx}                    % Formatação padronizada de unidades do SI (crucial para Engenharia)
\sisetup{output-decimal-marker = {,}}   % Configura a vírgula como separador decimal para PT-BR/PT-PT

% --- 3. LAYOUT, PÁGINA E ESPAÇAMENTO ---
\usepackage{geometry}       % Ajuste flexível das margens da página
\geometry{
    a4paper,
    top=30mm,
    bottom=20mm,
    left=30mm,
    right=20mm,
    % includehead,
    % includefoot,
    headheight=1.4cm,
    headsep=0.4cm,
    % headsep=0.6cm,
    footskip=1cm
    % footskip=1.4cm
}

\usepackage{fancyhdr}               % Customização avançada de cabeçalhos e rodapés
\usepackage{lastpage}               % Permite referenciar o número da última página (ex: "2 de 5")
\usepackage{setspace}               % Controle do espaçamento entre linhas (\onehalfspacing, etc.)
\usepackage{multicol}               % Suporte para layouts de múltiplas colunas na mesma página
\usepackage{indentfirst}            % Garante que o primeiro parágrafo de cada seção seja indentado

% --- 4. ELEMENTOS GRÁFICOS E CORES ---
\usepackage{graphicx}               % Inclusão de imagens e gráficos externos
\usepackage{transparent}            % Permite definir opacidade de imagens (ex.: logos com 70% de opacidade)
\usepackage{xcolor}                 % Suporte avançado e paletes de cores customizadas
\usepackage{tikz}                   % Ferramenta poderosa para criação de diagramas nativos
\usetikzlibrary{calc}               % Permite aritmética de coordenadas (usada na faixa do cabeçalho)
\usepackage{pgfplots}               % Geração de gráficos matemáticos de alta precisão
\pgfplotsset{compat=1.18}           % Garante compatibilidade do pgfplots

% --- 5. TABELAS PROFISSIONAIS ---
\usepackage{booktabs}               % Linhas de tabela padrão publicação (\toprule, \midrule, \bottomrule)
\usepackage{array}                  % Extensões para propriedades das colunas em tabelas
\usepackage{colortbl}               % Permite colorir linhas/colunas/células de tabelas (\rowcolor, \cellcolor)
\usepackage{multirow}               % Permite mesclar linhas em tabelas
\usepackage{tabularx}               % Tabelas com largura total fixa e colunas auto-ajustáveis (X)
\usepackage{longtable}              % Suporte para tabelas que quebram ao longo de múltiplas páginas

% --- 6. FIGURAS, LEGENDAS E FLUTUANTES ---
\usepackage{float}                  % Controle estrito de posicionamento (modificador [H])
\usepackage{caption}                % Customização estética das legendas de figuras e tabelas
\usepackage{subcaption}             % Criação de subfiguras e sub-legendas organizadas

% Legendas de figuras/tabelas usam o tamanho "menor" do padrão (Arial 10)
\DeclareCaptionFont{arialdez}{\fontsize{10}{12}\selectfont}
\captionsetup{font=arialdez, labelfont={arialdez,bf}}

% --- 7. LISTAS E ENUMERAÇÕES ---
\usepackage{enumitem}               % Customização total de listas (itemize, enumerate, description)

% --- 8. CÓDIGO-FONTE E CAIXAS DE DESTAQUE ---
\usepackage{listings}               % Inclusão e estilização de blocos de código-fonte
\usepackage{tcolorbox}              % Criação de caixas de texto coloridas e estilizadas de alta qualidade
\tcbuselibrary{skins,breakable}     % Permite caixas que quebram páginas e estilos complexos

% --- 9. TÍTULOS E SÍMBOLOS DIVERSOS ---
\usepackage{titlesec}               % Customização visual de \section, \subsection etc.
\usepackage{pifont}                 % Símbolos diversos (ex: \ding{51} para marcações de check)

% --- 10. REFERÊNCIAS E HIPERLINKS (Devem ser carregados por último) ---
\usepackage{url}                    % Formatação correta de URLs
\usepackage{hyperref}               % Transforma referências e sumários em links clicáveis interativos
\hypersetup{
    colorlinks=true,                % True para links coloridos, False para caixas coloridas
    linkcolor=black,                % Cor dos links internos (equações, seções)
    citecolor=black,                % Cor das citações bibliográficas
    urlcolor=black                  % Cor de links externos da internet
}
\usepackage{cleveref}               % Referências cruzadas inteligentes (escreve "Figura 1" automaticamente)

% Configuração de tradução do cleveref para português (garantia adicional)
\crefname{figure}{figura}{figuras}
\Crefname{figure}{Figura}{Figuras}
\crefname{table}{tabela}{tabelas}
\Crefname{table}{Tabela}{Tabelas}
\crefname{equation}{equação}{equações}
\Crefname{equation}{Equação}{Equações}

% --------------------------------------------------------------
% Aviso "Token not allowed in a PDF string (Unicode): removing
% `subscript'" do hyperref: aparece sempre que um campo de formulário
% (\TextField/\CheckBox, usados em campos-formulario.tex) é criado
% dentro de uma macro — é um comportamento interno conhecido do
% hyperref com AcroForm, inofensivo (não afeta o PDF gerado). A linha
% abaixo apenas filtra esse aviso específico da saída do compilador.
\usepackage{silence}
\WarningFilter{hyperref}{Token not allowed in a PDF string} antes de
% \begin{document}.
% ====================================================================

% --- 0. COMPILAÇÃO E FONTE ---
\usepackage[T1]{fontenc}                    % Codificação de fonte necessária para o pdfLaTeX
\usepackage[utf8]{inputenc}                 % Garante o suporte a acentuação nativa UTF-8
\usepackage{helvet}                         % Carrega a família Helvetica (equivalente visual ao Arial)
\renewcommand{\familydefault}{\sfdefault}   % Define a fonte sem serifa (Helvetica) como padrão do documento

% --- 1. IDIOMA E AJUSTES TIPOGRÁFICOS ---
\usepackage[brazil]{babel}          % Tradução de termos padrão (Figuras, Tabelas, Sumário)
\usepackage{microtype}              % Melhorias tipográficas estruturais (ajuste de espaçamento)

% --- 2. MATEMÁTICA AVANÇADA ---
\usepackage{amsmath}                    % Comandos matemáticos básicos da AMS
\usepackage{amsfonts}                   % Fontes matemáticas adicionais
\usepackage{amssymb}                    % Símbolos matemáticos diversos
\usepackage{mathtools}                  % Extensões e correções para o amsmath
\usepackage{bm}                         % Negrito para símbolos matemáticos (vetores, matrizes)
\usepackage{siunitx}                    % Formatação padronizada de unidades do SI (crucial para Engenharia)
\sisetup{output-decimal-marker = {,}}   % Configura a vírgula como separador decimal para PT-BR/PT-PT

% --- 3. LAYOUT, PÁGINA E ESPAÇAMENTO ---
\usepackage{geometry}       % Ajuste flexível das margens da página
\geometry{
    a4paper,
    top=30mm,
    bottom=20mm,
    left=30mm,
    right=20mm,
    % includehead,
    % includefoot,
    headheight=1.4cm,
    headsep=0.4cm,
    % headsep=0.6cm,
    footskip=1cm
    % footskip=1.4cm
}

\usepackage{fancyhdr}               % Customização avançada de cabeçalhos e rodapés
\usepackage{lastpage}               % Permite referenciar o número da última página (ex: "2 de 5")
\usepackage{setspace}               % Controle do espaçamento entre linhas (\onehalfspacing, etc.)
\usepackage{multicol}               % Suporte para layouts de múltiplas colunas na mesma página
\usepackage{indentfirst}            % Garante que o primeiro parágrafo de cada seção seja indentado

% --- 4. ELEMENTOS GRÁFICOS E CORES ---
\usepackage{graphicx}               % Inclusão de imagens e gráficos externos
\usepackage{transparent}            % Permite definir opacidade de imagens (ex.: logos com 70% de opacidade)
\usepackage{xcolor}                 % Suporte avançado e paletes de cores customizadas
\usepackage{tikz}                   % Ferramenta poderosa para criação de diagramas nativos
\usetikzlibrary{calc}               % Permite aritmética de coordenadas (usada na faixa do cabeçalho)
\usepackage{pgfplots}               % Geração de gráficos matemáticos de alta precisão
\pgfplotsset{compat=1.18}           % Garante compatibilidade do pgfplots

% --- 5. TABELAS PROFISSIONAIS ---
\usepackage{booktabs}               % Linhas de tabela padrão publicação (\toprule, \midrule, \bottomrule)
\usepackage{array}                  % Extensões para propriedades das colunas em tabelas
\usepackage{colortbl}               % Permite colorir linhas/colunas/células de tabelas (\rowcolor, \cellcolor)
\usepackage{multirow}               % Permite mesclar linhas em tabelas
\usepackage{tabularx}               % Tabelas com largura total fixa e colunas auto-ajustáveis (X)
\usepackage{longtable}              % Suporte para tabelas que quebram ao longo de múltiplas páginas

% --- 6. FIGURAS, LEGENDAS E FLUTUANTES ---
\usepackage{float}                  % Controle estrito de posicionamento (modificador [H])
\usepackage{caption}                % Customização estética das legendas de figuras e tabelas
\usepackage{subcaption}             % Criação de subfiguras e sub-legendas organizadas

% Legendas de figuras/tabelas usam o tamanho "menor" do padrão (Arial 10)
\DeclareCaptionFont{arialdez}{\fontsize{10}{12}\selectfont}
\captionsetup{font=arialdez, labelfont={arialdez,bf}}

% --- 7. LISTAS E ENUMERAÇÕES ---
\usepackage{enumitem}               % Customização total de listas (itemize, enumerate, description)

% --- 8. CÓDIGO-FONTE E CAIXAS DE DESTAQUE ---
\usepackage{listings}               % Inclusão e estilização de blocos de código-fonte
\usepackage{tcolorbox}              % Criação de caixas de texto coloridas e estilizadas de alta qualidade
\tcbuselibrary{skins,breakable}     % Permite caixas que quebram páginas e estilos complexos

% --- 9. TÍTULOS E SÍMBOLOS DIVERSOS ---
\usepackage{titlesec}               % Customização visual de \section, \subsection etc.
\usepackage{pifont}                 % Símbolos diversos (ex: \ding{51} para marcações de check)

% --- 10. REFERÊNCIAS E HIPERLINKS (Devem ser carregados por último) ---
\usepackage{url}                    % Formatação correta de URLs
\usepackage{hyperref}               % Transforma referências e sumários em links clicáveis interativos
\hypersetup{
    colorlinks=true,                % True para links coloridos, False para caixas coloridas
    linkcolor=black,                % Cor dos links internos (equações, seções)
    citecolor=black,                % Cor das citações bibliográficas
    urlcolor=black                  % Cor de links externos da internet
}
\usepackage{cleveref}               % Referências cruzadas inteligentes (escreve "Figura 1" automaticamente)

% Configuração de tradução do cleveref para português (garantia adicional)
\crefname{figure}{figura}{figuras}
\Crefname{figure}{Figura}{Figuras}
\crefname{table}{tabela}{tabelas}
\Crefname{table}{Tabela}{Tabelas}
\crefname{equation}{equação}{equações}
\Crefname{equation}{Equação}{Equações}

% --------------------------------------------------------------
% Aviso "Token not allowed in a PDF string (Unicode): removing
% `subscript'" do hyperref: aparece sempre que um campo de formulário
% (\TextField/\CheckBox, usados em campos-formulario.tex) é criado
% dentro de uma macro — é um comportamento interno conhecido do
% hyperref com AcroForm, inofensivo (não afeta o PDF gerado). A linha
% abaixo apenas filtra esse aviso específico da saída do compilador.
\usepackage{silence}
\WarningFilter{hyperref}{Token not allowed in a PDF string} antes de
% \begin{document}.
% ====================================================================

% --- 0. COMPILAÇÃO E FONTE ---
\usepackage[T1]{fontenc}                    % Codificação de fonte necessária para o pdfLaTeX
\usepackage[utf8]{inputenc}                 % Garante o suporte a acentuação nativa UTF-8
\usepackage{helvet}                         % Carrega a família Helvetica (equivalente visual ao Arial)
\renewcommand{\familydefault}{\sfdefault}   % Define a fonte sem serifa (Helvetica) como padrão do documento

% --- 1. IDIOMA E AJUSTES TIPOGRÁFICOS ---
\usepackage[brazil]{babel}          % Tradução de termos padrão (Figuras, Tabelas, Sumário)
\usepackage{microtype}              % Melhorias tipográficas estruturais (ajuste de espaçamento)

% --- 2. MATEMÁTICA AVANÇADA ---
\usepackage{amsmath}                    % Comandos matemáticos básicos da AMS
\usepackage{amsfonts}                   % Fontes matemáticas adicionais
\usepackage{amssymb}                    % Símbolos matemáticos diversos
\usepackage{mathtools}                  % Extensões e correções para o amsmath
\usepackage{bm}                         % Negrito para símbolos matemáticos (vetores, matrizes)
\usepackage{siunitx}                    % Formatação padronizada de unidades do SI (crucial para Engenharia)
\sisetup{output-decimal-marker = {,}}   % Configura a vírgula como separador decimal para PT-BR/PT-PT

% --- 3. LAYOUT, PÁGINA E ESPAÇAMENTO ---
\usepackage{geometry}       % Ajuste flexível das margens da página
\geometry{
    a4paper,
    top=30mm,
    bottom=20mm,
    left=30mm,
    right=20mm,
    % includehead,
    % includefoot,
    headheight=1.4cm,
    headsep=0.4cm,
    % headsep=0.6cm,
    footskip=1cm
    % footskip=1.4cm
}

\usepackage{fancyhdr}               % Customização avançada de cabeçalhos e rodapés
\usepackage{lastpage}               % Permite referenciar o número da última página (ex: "2 de 5")
\usepackage{setspace}               % Controle do espaçamento entre linhas (\onehalfspacing, etc.)
\usepackage{multicol}               % Suporte para layouts de múltiplas colunas na mesma página
\usepackage{indentfirst}            % Garante que o primeiro parágrafo de cada seção seja indentado

% --- 4. ELEMENTOS GRÁFICOS E CORES ---
\usepackage{graphicx}               % Inclusão de imagens e gráficos externos
\usepackage{transparent}            % Permite definir opacidade de imagens (ex.: logos com 70% de opacidade)
\usepackage{xcolor}                 % Suporte avançado e paletes de cores customizadas
\usepackage{tikz}                   % Ferramenta poderosa para criação de diagramas nativos
\usetikzlibrary{calc}               % Permite aritmética de coordenadas (usada na faixa do cabeçalho)
\usepackage{pgfplots}               % Geração de gráficos matemáticos de alta precisão
\pgfplotsset{compat=1.18}           % Garante compatibilidade do pgfplots

% --- 5. TABELAS PROFISSIONAIS ---
\usepackage{booktabs}               % Linhas de tabela padrão publicação (\toprule, \midrule, \bottomrule)
\usepackage{array}                  % Extensões para propriedades das colunas em tabelas
\usepackage{colortbl}               % Permite colorir linhas/colunas/células de tabelas (\rowcolor, \cellcolor)
\usepackage{multirow}               % Permite mesclar linhas em tabelas
\usepackage{tabularx}               % Tabelas com largura total fixa e colunas auto-ajustáveis (X)
\usepackage{longtable}              % Suporte para tabelas que quebram ao longo de múltiplas páginas

% --- 6. FIGURAS, LEGENDAS E FLUTUANTES ---
\usepackage{float}                  % Controle estrito de posicionamento (modificador [H])
\usepackage{caption}                % Customização estética das legendas de figuras e tabelas
\usepackage{subcaption}             % Criação de subfiguras e sub-legendas organizadas

% Legendas de figuras/tabelas usam o tamanho "menor" do padrão (Arial 10)
\DeclareCaptionFont{arialdez}{\fontsize{10}{12}\selectfont}
\captionsetup{font=arialdez, labelfont={arialdez,bf}}

% --- 7. LISTAS E ENUMERAÇÕES ---
\usepackage{enumitem}               % Customização total de listas (itemize, enumerate, description)

% --- 8. CÓDIGO-FONTE E CAIXAS DE DESTAQUE ---
\usepackage{listings}               % Inclusão e estilização de blocos de código-fonte
\usepackage{tcolorbox}              % Criação de caixas de texto coloridas e estilizadas de alta qualidade
\tcbuselibrary{skins,breakable}     % Permite caixas que quebram páginas e estilos complexos

% --- 9. TÍTULOS E SÍMBOLOS DIVERSOS ---
\usepackage{titlesec}               % Customização visual de \section, \subsection etc.
\usepackage{pifont}                 % Símbolos diversos (ex: \ding{51} para marcações de check)

% --- 10. REFERÊNCIAS E HIPERLINKS (Devem ser carregados por último) ---
\usepackage{url}                    % Formatação correta de URLs
\usepackage{hyperref}               % Transforma referências e sumários em links clicáveis interativos
\hypersetup{
    colorlinks=true,                % True para links coloridos, False para caixas coloridas
    linkcolor=black,                % Cor dos links internos (equações, seções)
    citecolor=black,                % Cor das citações bibliográficas
    urlcolor=black                  % Cor de links externos da internet
}
\usepackage{cleveref}               % Referências cruzadas inteligentes (escreve "Figura 1" automaticamente)

% Configuração de tradução do cleveref para português (garantia adicional)
\crefname{figure}{figura}{figuras}
\Crefname{figure}{Figura}{Figuras}
\crefname{table}{tabela}{tabelas}
\Crefname{table}{Tabela}{Tabelas}
\crefname{equation}{equação}{equações}
\Crefname{equation}{Equação}{Equações}

% --------------------------------------------------------------
% Aviso "Token not allowed in a PDF string (Unicode): removing
% `subscript'" do hyperref: aparece sempre que um campo de formulário
% (\TextField/\CheckBox, usados em campos-formulario.tex) é criado
% dentro de uma macro — é um comportamento interno conhecido do
% hyperref com AcroForm, inofensivo (não afeta o PDF gerado). A linha
% abaixo apenas filtra esse aviso específico da saída do compilador.
\usepackage{silence}
\WarningFilter{hyperref}{Token not allowed in a PDF string}